% ================================================================ TD 1
\td{1}{Dockerisation de My Favorite Places}

My Favorite Places est composée d'un client React, d'une API Node.js/Express et d'une base PostgreSQL. Dans ce TD, j'ai mis ces trois parties dans des conteneurs, puis je les ai orchestrées avec Docker Compose.

% ---------------------------------------------------------------- 1
\partie{1}{Test initial avec Bruno}

\subsection{Lancer une base PostgreSQL}

Le serveur ne démarre pas sans base. J'ai lancé PostgreSQL dans un conteneur avec le mot de passe attendu par \code{server/src/datasource.ts}.

\begin{commandes}
docker run -d --name mfp-db \
  -e POSTGRES_PASSWORD=supersecret \
  -p 5432:5432 \
  postgres:17
docker ps
\end{commandes}

\subsection{Lancer l'API en mode développement}

\begin{commandes}
cd server
yarn install
yarn dev
\end{commandes}

La console affiche \code{connected to PgSQL db} puis \code{server started on port 3000}.

\subsection{Tests avec la collection Bruno}

J'ai ouvert le dossier \code{mfp-bruno-api-test-collection/} dans Bruno et lancé les requêtes dans l'ordre, avec le compte \code{ahmedtest@gmail.com} et le mot de passe \code{azerty123}.

\subsubsection{Créer un utilisateur}

\capture[0.20\textheight]{td1-01-create-user.png}{Requête « Create user » sur \code{POST /api/users} : l'API répond \code{200 OK} et renvoie l'utilisateur créé, avec son mot de passe haché.}

\subsubsection{Se connecter}

La requête « Login (get token) » sur \code{POST /api/users/tokens} renvoie \code{200 OK} et un token. La même requête est montrée en 2.4, avec le serveur dans Docker.

\subsubsection{Consulter mon profil}

\capture[0.24\textheight]{td1-03-my-profile.png}{Requête « My profile » sur \code{GET /api/users/me}, avec le token dans l'en-tête \code{Authorization}. L'API répond \code{200 OK} et renvoie l'utilisateur connecté.}

% ---------------------------------------------------------------- 2
\partie{2}{Dockerisation du serveur}

\subsection{Le Dockerfile du serveur}

\begin{fichier}[docker]{server/Dockerfile}
FROM node:24-alpine3.24
WORKDIR /app
COPY package.json yarn.lock ./
RUN yarn install --frozen-lockfile
COPY tsconfig.json ./
COPY src src
RUN yarn build
EXPOSE 3000
CMD ["yarn", "start"]
\end{fichier}

J'ai aussi ajouté un \code{server/.dockerignore} qui écarte \code{node_modules} et \code{dist}. Les captures de ce TD ont été prises avec l'image \code{node:22}, que j'ai remplacée ensuite par \code{node:24-alpine3.24}.

\subsection{Adapter la connexion à la base}

Dans un conteneur, \code{localhost} désigne le conteneur lui-même. J'ai donc fait lire les paramètres de connexion dans des variables d'environnement.

\begin{fichier}[typescript]{server/src/datasource.ts}
host: process.env.DB_HOST || "localhost",
username: process.env.DB_USER || "postgres",
password: process.env.DB_PASSWORD || "supersecret",
database: process.env.DB_NAME || "postgres",
\end{fichier}

\subsection{Lancer la base et le serveur}

J'ai créé un réseau commun, sur lequel le serveur trouve la base sous le nom \code{mfp-db}.

\begin{commandes}[Commandes]
# 1. créer un réseau commun
docker network create mfp-net
# 2. lancer la base sur ce réseau
docker run -d --name mfp-db --network mfp-net \
  -e POSTGRES_PASSWORD=supersecret \
  postgres:17
# 3. construire l'image du serveur
docker build -t mfp-server ./server
# 4. lancer le serveur sur le même réseau
docker run -d --name mfp-server --network mfp-net \
  -e DB_HOST=mfp-db \
  -p 3000:3000 \
  mfp-server
\end{commandes}

\probleme{La commande 4 a échoué avec \code{address already in use} sur le port 3000, car le \code{yarn dev} de la partie 1 tournait encore.}{J'ai arrêté \code{yarn dev} avec Ctrl+C, supprimé le conteneur raté et relancé la commande 4.}

\begin{commandes}
docker rm -f mfp-server
# 4. lancer le serveur sur le même réseau
docker run -d --name mfp-server --network mfp-net \
  -e DB_HOST=mfp-db \
  -p 3000:3000 \
  mfp-server
# 5. vérifier
docker ps
docker logs mfp-server
\end{commandes}

\capture[0.34\textheight]{td1-05-server-docker-logs.png}{Les conteneurs \code{mfp-server} et \code{mfp-db} sont \code{Up}. Dans les logs, le serveur crée les tables \code{user} et \code{address}, puis se connecte à la base.}

\subsection{Test avec Bruno}

Grâce à \code{-p 3000:3000}, l'API répond à la même adresse qu'avec \code{yarn dev}. Comme la base était neuve, j'ai d'abord recréé l'utilisateur avec « Create user ».

\capture[0.20\textheight]{td1-06-login-docker.png}{Requête « Login (get token) » vers le serveur lancé dans Docker : \code{200 OK} et un token. Le serveur retrouve bien l'utilisateur dans \code{mfp-db}.}

% ---------------------------------------------------------------- 3
\partie{3}{Dockerisation du client}

\subsection{Le Dockerfile du client}

\begin{fichier}[docker]{client/Dockerfile}
FROM node:24-alpine3.24
WORKDIR /app
COPY package*.json yarn.lock ./
RUN yarn install --frozen-lockfile
COPY tsconfig*.json ./
COPY vite.config.ts ./
COPY index.html ./
COPY public public
COPY src src
EXPOSE 5173
CMD ["yarn", "dev", "--host"]
\end{fichier}

Sans l'option \code{--host}, Vite n'accepte que les connexions venant de l'intérieur du conteneur.

\subsection{Adapter l'adresse de l'API}

Vite fait suivre les appels \code{/api} au serveur. Comme pour la base, j'ai mis l'adresse du serveur dans une variable d'environnement.

\begin{fichier}[typescript]{client/vite.config.ts}
target: process.env.API_URL || "http://localhost:3000",
\end{fichier}

\subsection{Lancer le client}

La base et le serveur de la partie 2 tournent toujours.

\begin{commandes}[Commandes]
# 1. construire l'image du client
docker build -t mfp-client ./client
# 2. lancer le client sur le même réseau que le serveur
docker run -d --name mfp-client --network mfp-net \
  -e API_URL=http://mfp-server:3000 \
  -p 5173:5173 \
  mfp-client
# 3. vérifier
docker ps
docker logs mfp-client
\end{commandes}

\capture[0.34\textheight]{td1-07-client-docker.png}{Les trois conteneurs \code{mfp-client}, \code{mfp-server} et \code{mfp-db} sont \code{Up}. Les logs du client montrent Vite prêt sur le port 5173.}

\subsection{Test dans le navigateur}

\capture[0.20\textheight]{td1-08-client-browser.png}{Connexion sur \code{http://localhost:5173} avec le compte de test : le tableau de bord s'affiche.}

% ---------------------------------------------------------------- 4
\needspace{8cm}
\partie{4}{Orchestration avec Docker Compose}

\subsection{Le fichier compose.yml}

Chaque service remplace une commande \code{docker run} que je tapais à la main. Compose crée lui-même le réseau, sur lequel chaque service est joignable par son nom.

\begin{fichier}{compose.yml (extrait)}
services:
  db:
    image: postgres:17
    restart: unless-stopped
    environment:
      POSTGRES_USER: postgres
      POSTGRES_PASSWORD: supersecret
    ports:
      - "5432:5432"
    volumes:
      - db-data:/var/lib/postgresql/data
    healthcheck:
      test: ["CMD-SHELL", "pg_isready -U postgres"]
      interval: 2s
      timeout: 5s
      retries: 15

  server:
    build: ./server
    restart: unless-stopped
    ports:
      - "3000:3000"
    depends_on:
      db:
        condition: service_healthy
\end{fichier}

Le service \code{client} publie le port 5173 et dépend de \code{server}. Le volume \code{db-data} garde les données de la base quand le conteneur est recréé.

\subsection{Lancer l'application}

J'ai d'abord supprimé les conteneurs et le réseau créés à la main, qui occupaient encore les ports 3000 et 5173.

\begin{commandes}
docker rm -f mfp-client mfp-server mfp-db
docker network rm mfp-net

docker compose up -d --build
docker compose ps
docker compose logs server
\end{commandes}

\probleme{Avec un simple \code{depends_on}, le serveur s'arrêtait sur \code{ECONNREFUSED}, car PostgreSQL n'était pas encore prêt au démarrage du serveur.}{J'ai ajouté le \code{healthcheck} sur \code{db} et \code{condition: service_healthy} sur \code{server}.}

\capture{td1-09-compose-up.png}{Compose attend que \code{mfp-db-1} soit \code{Healthy} avant de démarrer le serveur et le client. Les trois services sont \code{Up} et le serveur se connecte à la base.}

\subsection{Test}

Le volume étant neuf, j'ai recréé l'utilisateur avec Bruno, puis je me suis connecté au client avec le même résultat qu'en 3.4. Pour tout arrêter sans perdre le volume \code{db-data} :

\begin{commandes}
docker compose down
\end{commandes}

\subsection{Simplification de la connexion à la base}

Avec Compose, la base s'appelle toujours \code{db}. J'ai donc retiré les variables d'environnement et écrit l'adresse directement dans le code.

\begin{fichier}[typescript]{server/src/datasource.ts}
host: "db",
username: "postgres",
password: "supersecret",
database: "postgres",
\end{fichier}

En contrepartie, le serveur ne trouve plus la base en dehors de Compose.

% ---------------------------------------------------------------- 5
% le titre est suivi du schéma (insécable) : place pour les deux
\Needspace{6cm}
\partie{5}{Communication entre les services}

\begin{center}
\begin{tikzpicture}[font=\small, >=stealth,
    boite/.style={draw, minimum height=7mm, inner sep=4pt},
    etiquette/.style={font=\footnotesize, inner sep=2pt}]
  \draw[dashed] (-1.4,-1.25) rectangle (1.4,1.35);
  \node[etiquette, anchor=south] at (0,1.35) {Machine (WSL)};
  \draw[dashed] (3.0,-1.25) rectangle (14.4,1.35);
  \node[etiquette, anchor=south] at (8.7,1.35) {Réseau Docker \code{mfp_default}};
  \node[boite] (nav) at (0,0.6) {Navigateur};
  \node[boite] (bruno) at (0,-0.6) {Bruno};
  \node[boite] (client) at (5.0,0.6) {\code{client} (nginx)};
  \node[boite] (server) at (9.4,0.6) {\code{server}};
  \node[boite] (db) at (13.0,0.6) {\code{db}};
  \draw[->] (nav) -- node[etiquette, above] {:5173} (client);
  \draw[->] (client) -- node[etiquette, above] {\code{/api}, :3000} (server);
  \draw[->] (server) -- node[etiquette, above] {:5432} (db);
  \draw[->] (bruno) -| node[etiquette, above, pos=0.3] {:3000} (server);
\end{tikzpicture}
\end{center}

\subsection{Entre la base et le serveur}

Les deux conteneurs sont sur le réseau \code{mfp_default} créé par Compose. Le serveur appelle la base par son nom de service, \code{db}, sur le port 5432, et Docker traduit ce nom en adresse IP. Le port 5432 n'est publié que pour consulter la base depuis ma machine.

\subsection{Entre le client et le serveur}

Le code React s'exécute dans le navigateur, hors du réseau Docker, et appelle donc l'API à l'adresse du site (\code{http://localhost:5173/api/...}). Le conteneur \code{client} relaie ces appels vers le service \code{server}, d'abord avec Vite puis avec nginx. Bruno, lui, appelle directement le serveur sur le port 3000 publié.

% ---------------------------------------------------------------- 6
\partie{6}{Démarrage, redémarrage et images légères}

\subsection{Démarrer le serveur quand la base est prête}

C'était déjà réglé en 4.2 avec le \code{healthcheck} sur \code{db} et \code{condition: service_healthy} sur \code{server}.

\subsection{Redémarrage automatique}

J'ai ajouté cette ligne à chaque service de \code{compose.yml}.

\begin{fichier}{compose.yml}
restart: unless-stopped
\end{fichier}

Pour le vérifier, j'ai simulé un plantage du serveur, puis un arrêt manuel.

\begin{commandes}
docker compose exec server pkill -f "node dist/index.js"
docker inspect mfp-server-1 --format 'status={{.State.Status}} restarts={{.RestartCount}}'
# status=running restarts=1

docker compose stop server
docker inspect mfp-server-1 --format 'status={{.State.Status}} restarts={{.RestartCount}}'
# status=exited restarts=1   (arrêt manuel : pas de redémarrage)

docker compose start server
\end{commandes}

\needspace{7cm}
\subsection{Alléger les images}

J'ai découpé chaque Dockerfile en trois étapes : \code{deps} installe les dépendances, \code{builder} construit l'application et \code{main} ne garde que ce qui sert à la lancer.

\subsubsection{Serveur}

L'image finale ne garde que le JavaScript compilé et les dépendances de production.

\begin{fichier}[docker]{server/Dockerfile}
FROM node:24-alpine3.24 AS deps
WORKDIR /app
COPY package.json yarn.lock ./
RUN yarn install --frozen-lockfile

FROM node:24-alpine3.24 AS builder
WORKDIR /app
COPY --from=deps /app/node_modules ./node_modules
COPY package.json ./
COPY src ./src
COPY tsconfig.json ./

RUN yarn build

FROM node:24-alpine3.24 AS main
WORKDIR /app

COPY --from=builder /app/dist ./dist
COPY package.json yarn.lock ./
RUN yarn install --frozen-lockfile --production && yarn cache clean

EXPOSE 3000

CMD ["yarn", "start"]
\end{fichier}

\subsubsection{Client}

Une fois construit, le client n'est plus qu'un ensemble de fichiers statiques. Les étapes \code{deps} et \code{builder} sont les mêmes que pour le serveur, mais l'image finale part de nginx.

\begin{fichier}[docker]{client/Dockerfile (étape finale)}
# ---- stage 3: production image, nginx serves the static files ----
FROM nginx:1.29-alpine AS main

COPY nginx.conf /etc/nginx/conf.d/default.conf
COPY --from=builder /app/dist /usr/share/nginx/html

EXPOSE 5173
\end{fichier}

Dans cette image, nginx reprend le rôle de Vite : il fait suivre les appels \code{/api/} au service \code{server}.

\begin{fichier}[nginx]{client/nginx.conf}
server {
    listen 5173;
    root /usr/share/nginx/html;

    # React routes (/signin, /dashboard...) all load index.html
    location / {
        try_files $uri $uri/ /index.html;
    }

    # API calls are forwarded to the server service
    location /api/ {
        proxy_pass http://server:3000;
    }
}
\end{fichier}

\subsubsection{Résultat}

\begin{commandes}
docker compose up -d --build
docker images
\end{commandes}

\begin{center}
\begin{tabular}{lrr}
  \toprule
  Image & Avant & Après \\
  \midrule
  \code{mfp-server} & 477 Mo & 337 Mo \\
  \code{mfp-client} & 1,49 Go & 93 Mo \\
  \bottomrule
\end{tabular}
\end{center}
